\documentclass[12pt,a4paper]{article}
\usepackage[utf8]{inputenc}
\usepackage[spanish]{babel}
\usepackage[T1]{fontenc}
\usepackage{geometry}
\usepackage{graphicx}
\usepackage{xcolor}
\usepackage{titlesec}
\usepackage{enumitem}
\usepackage{fancyhdr}
\usepackage{tcolorbox}
\usepackage{fontawesome5}
\usepackage{hyperref}
\usepackage{booktabs}
\usepackage{array}
\usepackage{mdframed}

\geometry{top=2.5cm, bottom=2.5cm, left=2.5cm, right=2.5cm}

% ── Colores corporativos ──────────────────────────────────────────
\definecolor{primario}{HTML}{C9A96E}
\definecolor{oscuro}{HTML}{1C1C1E}
\definecolor{gris}{HTML}{6B6B6B}
\definecolor{fondo}{HTML}{FAF7F2}
\definecolor{alerta}{HTML}{E74C3C}
\definecolor{exito}{HTML}{2ECC71}
\definecolor{info}{HTML}{3498DB}

% ── Secciones estilo ────────────────────────────────────────────
\titleformat{\section}{\Large\bfseries\color{oscuro}}{}{0em}{}[\color{primario}\titlerule]
\titleformat{\subsection}{\large\bfseries\color{oscuro}}{}{0em}{\faAngleRight\ }
\titleformat{\subsubsection}{\normalsize\bfseries\color{gris}}{}{0em}{}

% ── Cajas de información ────────────────────────────────────────
\tcbuselibrary{skins,breakable}
\newtcolorbox{notabox}{
  colback=fondo, colframe=primario, coltitle=white, fonttitle=\bfseries,
  title={\faInfoCircle\ Nota}, breakable, arc=4pt
}
\newtcolorbox{alertabox}{
  colback=red!5, colframe=alerta, coltitle=white, fonttitle=\bfseries,
  title={\faExclamationTriangle\ Importante}, breakable, arc=4pt
}
\newtcolorbox{pasobox}[1]{
  colback=blue!3, colframe=info, coltitle=white, fonttitle=\bfseries,
  title={\faCheckCircle\ #1}, breakable, arc=4pt
}

% ── Encabezado/pie ──────────────────────────────────────────────
\pagestyle{fancy}
\fancyhf{}
\fancyhead[L]{\textcolor{primario}{\textbf{MUYMUY Beauty Studio}}}
\fancyhead[R]{\textcolor{gris}{\small Manual de Empleado}}
\fancyfoot[C]{\textcolor{gris}{\small Página \thepage}}
\renewcommand{\headrulewidth}{0.5pt}
\renewcommand{\headrule}{\hbox to\headwidth{\color{primario}\leaders\hrule height \headrulewidth\hfill}}

% ── Hyperlinks ───────────────────────────────────────────────────
\hypersetup{colorlinks=true, linkcolor=primario, urlcolor=primario}

% ════════════════════════════════════════════════════════════════
\begin{document}

% ── Portada ─────────────────────────────────────────────────────
\begin{titlepage}
  \centering
  \vspace*{4cm}
  {\color{primario}\rule{\linewidth}{1pt}}\\[1cm]
  {\color{oscuro}\Huge\bfseries MUYMUY Beauty Studio}\\[0.5cm]
  {\color{primario}\Large Manual de Usuario}\\[0.3cm]
  {\color{oscuro}\LARGE\bfseries ROL: EMPLEADO}\\[1cm]
  {\color{primario}\rule{\linewidth}{1pt}}\\[2cm]
  {\color{gris}\large Versión 1.1 --- \today}\\[0.5cm]
  {\color{gris}\normalsize Sistema Administrativo Interno}
  \vfill
  {\color{oscuro}\small Documento confidencial --- Solo para uso interno}
\end{titlepage}
\nopagecolor

% ── Índice ──────────────────────────────────────────────────────
\tableofcontents
\newpage

% ════════════════════════════════════════════════════════════════
\section{Introducción}

Este manual describe el uso del sistema administrativo de \textbf{MUYMUY Beauty Studio} para empleados. Como empleado, tienes acceso a las secciones necesarias para operar tu turno diario de forma completa y autónoma.

\begin{notabox}
Tu usuario solo da acceso a: \textbf{Agenda, Clientes, Asistencia, Venta Directa, Caja} y \textbf{Vacaciones}. Las secciones de análisis, inventario y configuración son exclusivas del administrador.
\end{notabox}

\subsection{Acceso al Sistema}

\begin{pasobox}{Cómo iniciar sesión}
\begin{enumerate}[leftmargin=*]
  \item Abre el navegador y entra a la URL del sistema.
  \item Escribe tu correo electrónico y contraseña.
  \item Haz clic en \textbf{Iniciar sesión}.
  \item El sistema te llevará directamente a la \textbf{Agenda}.
\end{enumerate}
\end{pasobox}

\begin{alertabox}
Nunca compartas tu contraseña con otras personas. Cada empleada tiene su propio acceso único. Al terminar tu turno, usa el botón \textbf{Cerrar sesión} en la parte inferior del menú lateral.
\end{alertabox}

\subsection{Navegación General}

Al entrar al sistema verás una barra lateral izquierda con las secciones disponibles. En la parte superior de esa barra aparece el \textbf{selector de sucursal}~--- asegúrate de que esté seleccionada tu sucursal antes de comenzar a trabajar.

% ════════════════════════════════════════════════════════════════
\newpage
\section{Asistencia}

La sección de Asistencia es lo primero que debes usar al llegar al trabajo. Registra tu llegada, salida a comer y regreso para que el sistema sepa que estás disponible en la agenda.

\subsection{Registrar Entrada}

\begin{pasobox}{Al llegar al trabajo}
\begin{enumerate}[leftmargin=*]
  \item Ve a la sección \textbf{Asistencia} en el menú lateral.
  \item Busca tu nombre en la lista de empleadas de tu sucursal.
  \item Haz clic en el botón \textbf{Registrar Entrada}.
  \item El sistema mostrará la hora de entrada y tu nombre aparecerá en la agenda.
\end{enumerate}
\end{pasobox}

\begin{notabox}
Una vez que registras tu \textbf{Entrada}, tu columna en la Agenda se activa y puedes recibir citas. Si no registras entrada, aparecerás como no disponible.
\end{notabox}

\subsection{Ir a Comer / Regreso de Comida}

Cuando salgas a comer, es importante registrarlo porque el sistema \textbf{bloquea automáticamente} tu horario en la agenda durante ese tiempo.

\begin{pasobox}{Para registrar salida a comer}
\begin{enumerate}[leftmargin=*]
  \item Ve a \textbf{Asistencia}.
  \item Haz clic en \textbf{Ir a Comer} junto a tu nombre.
  \item Tu columna en la Agenda mostrará un bloqueo automático.
  \item Al regresar, haz clic en \textbf{Regreso de Comida}.
  \item El bloqueo se elimina y vuelves a estar disponible.
\end{enumerate}
\end{pasobox}

\begin{alertabox}
Este es el \textbf{único mecanismo} mediante el cual una empleada puede bloquear su propio horario. Los empleados \textbf{no tienen acceso} al botón de ``Bloquear Horario'' manual --- esa función es exclusiva de administradores.
\end{alertabox}

\subsection{Registrar Salida}

Al terminar tu turno:

\begin{pasobox}{Al terminar el turno}
\begin{enumerate}[leftmargin=*]
  \item Ve a \textbf{Asistencia}.
  \item Haz clic en \textbf{Registrar Salida} junto a tu nombre.
  \item El sistema guarda tu hora de salida y cierra tu jornada.
\end{enumerate}
\end{pasobox}

% ════════════════════════════════════════════════════════════════
\newpage
\section{Agenda}

La Agenda es el centro de operaciones del sistema. Aquí puedes ver todas las citas de la semana, agendar nuevas y gestionar el estado de cada cita.

\subsection{Vista General de la Agenda}

La agenda muestra una cuadrícula semanal con:
\begin{itemize}[leftmargin=*]
  \item \textbf{Columnas:} Una por cada día de la semana (Lun--Dom).
  \item \textbf{Filas:} Bloques de tiempo desde la mañana hasta la noche.
  \item \textbf{Empleadas:} Cada empleada tiene su propia sección por columna.
  \item \textbf{Colores:} Las citas muestran colores según su estado (programada, en curso, finalizada, cancelada).
\end{itemize}

Puedes navegar entre semanas con las flechas \textbf{◄ ►} o volver a la semana actual con el botón \textbf{Hoy}.

\subsection{Agendar una Nueva Cita}

\begin{pasobox}{Cómo crear una cita}
\begin{enumerate}[leftmargin=*]
  \item Haz clic en un slot (casilla) vacío en la cuadrícula de la agenda.
  \item Se abrirá el \textbf{buscador de clientes}. Escribe el nombre o teléfono del cliente.
  \item Si el cliente existe, selecciónalo de la lista.
  \item Si es un cliente nuevo, haz clic en \textbf{Nuevo Cliente} y llena el formulario.
  \item Una vez seleccionado el cliente, elige los servicios que va a recibir.
  \item Confirma la cita. Aparecerá en la cuadrícula.
\end{enumerate}
\end{pasobox}

\begin{notabox}
También puedes ir primero a \textbf{Clientes}, buscar al cliente y desde su perfil hacer clic en ``Agendar Cita'', que te llevará directamente a la agenda con ese cliente ya seleccionado.
\end{notabox}

\subsection{Gestionar una Cita Existente}

\begin{pasobox}{Cómo ver o modificar una cita}
\begin{enumerate}[leftmargin=*]
  \item Haz clic sobre cualquier cita en la cuadrícula.
  \item Se abre el panel de \textbf{Gestión de Cita} con todos los detalles.
  \item Desde ahí puedes:
  \begin{itemize}
    \item Ver los servicios agendados.
    \item Cambiar el estado (Programada, En curso, Finalizada, Cancelada, No asistió).
    \item Iniciar el proceso de \textbf{cobro} cuando la cita termina.
  \end{itemize}
\end{enumerate}
\end{pasobox}

\subsection{Proceso de Cobro de una Cita}

Cuando la clienta termina su servicio, debes realizar el cobro desde la agenda:

\begin{pasobox}{Cobrar una cita}
\begin{enumerate}[leftmargin=*]
  \item Haz clic en la cita finalizada.
  \item En el panel de gestión, haz clic en \textbf{Cobrar / Validar}.
  \item Se abre el flujo de checkout. Revisa los servicios y precios.
  \item Puedes agregar productos adicionales si la clienta compró algo.
  \item Selecciona el \textbf{método de pago} (Efectivo, Tarjeta, Transferencia, etc.).
  \item Haz clic en \textbf{Confirmar Pago}. Se genera el ticket automáticamente.
\end{enumerate}
\end{pasobox}

\subsection{Bloqueos en la Agenda}

\begin{alertabox}
Como empleada, \textbf{no puedes crear bloqueos manuales} en la agenda. Los únicos bloqueos que puedes generar son los automáticos al registrar \textbf{Ir a Comer} en la sección de Asistencia. Si necesitas bloquear un horario por otra razón, comunícalo al administrador.
\end{alertabox}

Si ves un bloqueo en la agenda (celda gris o con candado), puedes hacer clic en él para ver el motivo y la duración, pero no podrás eliminarlo.

% ════════════════════════════════════════════════════════════════
\newpage
\section{Clientes}

La sección de Clientes te permite buscar, registrar y consultar la información de las clientas del salón.

\subsection{Buscar un Cliente}

\begin{pasobox}{Búsqueda de cliente}
\begin{enumerate}[leftmargin=*]
  \item Ve a la sección \textbf{Clientes}.
  \item Usa la barra de búsqueda para escribir nombre, teléfono o correo.
  \item Los resultados aparecen en tiempo real conforme escribes.
  \item Haz clic en el cliente para ver su perfil completo.
\end{enumerate}
\end{pasobox}

\subsection{Registrar un Nuevo Cliente}

\begin{pasobox}{Alta de cliente nuevo}
\begin{enumerate}[leftmargin=*]
  \item En la sección \textbf{Clientes}, haz clic en \textbf{Nuevo Cliente}.
  \item Llena los campos requeridos:
    \begin{itemize}
      \item Nombre completo \textit{(obligatorio)}
      \item Teléfono \textit{(obligatorio)}
      \item Correo electrónico \textit{(opcional)}
      \item Fecha de nacimiento \textit{(opcional, para felicitaciones)}
      \item Notas \textit{(alergias, preferencias, etc.)}
    \end{itemize}
  \item Haz clic en \textbf{Guardar}.
\end{enumerate}
\end{pasobox}

\subsection{Perfil del Cliente}

Desde el perfil de un cliente puedes ver:
\begin{itemize}[leftmargin=*]
  \item Historial completo de visitas y servicios.
  \item Total gastado en el salón.
  \item Última visita.
  \item Botón directo para \textbf{Agendar una nueva cita}.
\end{itemize}

% ════════════════════════════════════════════════════════════════
\newpage
\section{Venta Directa}

La Venta Directa te permite registrar ventas de productos o servicios sin necesidad de una cita agendada previamente.

\subsection{¿Cuándo usar Venta Directa?}

Usa esta sección cuando:
\begin{itemize}[leftmargin=*]
  \item Un cliente llega a comprar un producto sin cita.
  \item Se realiza un servicio exprés sin cita previa.
  \item Necesitas hacer una venta rápida sin pasar por la agenda.
\end{itemize}

\subsection{Registrar una Venta}

\begin{pasobox}{Proceso de venta directa}
\begin{enumerate}[leftmargin=*]
  \item Ve a \textbf{Venta Directa}.
  \item Selecciona o busca al cliente (puede ser venta sin cliente registrado).
  \item Agrega los productos o servicios que se están vendiendo.
  \item El sistema calcula el total automáticamente.
  \item Selecciona el método de pago.
  \item Haz clic en \textbf{Confirmar Venta}. Se genera el ticket.
\end{enumerate}
\end{pasobox}

% ════════════════════════════════════════════════════════════════
\newpage
\section{Caja}

La sección de Caja te permite gestionar el dinero en efectivo del turno.

\subsection{Abrir Turno de Caja}

Al iniciar tu turno de caja (si eres responsable de ella):

\begin{pasobox}{Abrir caja}
\begin{enumerate}[leftmargin=*]
  \item Ve a la sección \textbf{Caja}.
  \item Haz clic en \textbf{Abrir Turno}.
  \item Ingresa el monto de efectivo con el que inicias (fondo de caja).
  \item Confirma la apertura.
\end{enumerate}
\end{pasobox}

\subsection{Registrar Movimientos de Caja}

Puedes registrar entradas o salidas de efectivo durante el turno:

\begin{pasobox}{Registrar movimiento}
\begin{enumerate}[leftmargin=*]
  \item En \textbf{Caja}, selecciona \textbf{Nuevo Movimiento}.
  \item Elige el tipo: \textbf{Entrada} o \textbf{Salida}.
  \item Selecciona el concepto del movimiento del menú desplegable.
  \item Ingresa el monto.
  \item Agrega una nota si es necesario.
  \item Confirma el movimiento.
\end{enumerate}
\end{pasobox}

\subsection{Cerrar Turno de Caja}

Al terminar el turno:

\begin{pasobox}{Cerrar caja}
\begin{enumerate}[leftmargin=*]
  \item Ve a \textbf{Caja} y haz clic en \textbf{Cerrar Turno}.
  \item Cuenta el efectivo físico que tienes y escribe el monto.
  \item El sistema calculará automáticamente si hay diferencia (faltante o sobrante).
  \item Confirma el cierre.
\end{enumerate}
\end{pasobox}

% ════════════════════════════════════════════════════════════════
\newpage
\section{Vacaciones}

El sistema cuenta con un módulo para gestionar tus solicitudes de descanso de forma digital y organizada.

\subsection{Solicitar Vacaciones}

\begin{pasobox}{Nueva Solicitud}
\begin{enumerate}[leftmargin=*]
  \item Ve a la sección \textbf{Vacaciones} en el menú lateral.
  \item Haz clic en el botón \textbf{Nueva Solicitud}.
  \item Selecciona la \textbf{Fecha de Inicio} y la \textbf{Fecha de Fin}.
  \item Escribe una breve nota o motivo si es necesario.
  \item Haz clic en \textbf{Enviar Solicitud}.
\end{enumerate}
\end{pasobox}

\subsection{Estado de Solicitudes}

En la pantalla principal de Vacaciones verás el listado de tus solicitudes enviadas y su estado:
\begin{itemize}[leftmargin=*]
  \item \textbf{Pendiente:} El administrador aún no ha revisado la solicitud.
  \item \textbf{Aprobada:} Tus vacaciones han sido autorizadas. El sistema \textbf{bloqueará automáticamente} tu horario en la agenda para esos días.
  \item \textbf{Rechazada:} La solicitud no pudo ser autorizada. Podrás ver el motivo en la nota del administrador.
\end{itemize}

\subsection{Extensiones de Vacaciones}

Si ya tienes una solicitud aprobada y necesitas agregar más días:

\begin{pasobox}{Solicitar Extensión}
\begin{enumerate}[leftmargin=*]
  \item Busca tu solicitud aprobada en el listado.
  \item Haz clic en el botón \textbf{Solicitar Extensión}.
  \item Selecciona las nuevas fechas y envía para aprobación.
\end{enumerate}
\end{pasobox}

\begin{alertabox}
Una vez que tus vacaciones son \textbf{Aprobadas}, no es necesario que el administrador bloquee tu agenda manualmente; el sistema lo hace por ti para asegurar que no recibas citas en esos días.
\end{alertabox}

% ════════════════════════════════════════════════════════════════
\newpage
\section{Preguntas Frecuentes}

\begin{description}[leftmargin=0pt]
  \item[\textbf{¿Qué hago si no puedo iniciar sesión?}]
    Verifica que el correo y la contraseña sean correctos. Si olvidaste tu contraseña, pide al administrador que la restablezca.

  \item[\textbf{¿Por qué no aparezco en la agenda?}]
    Debes registrar tu \textbf{Entrada} en la sección de Asistencia para que tu columna se active en la agenda.

  \item[\textbf{¿Puedo cambiar la sucursal activa?}]
    Sí, desde el selector en la parte superior de la barra lateral. Asegúrate de seleccionar tu sucursal correcta.

  \item[\textbf{¿Qué hago si cometí un error en una venta?}]
    No puedes eliminar tickets una vez confirmados. Notifica al administrador para que haga la corrección correspondiente.

  \item[\textbf{¿Por qué no veo los botones de bloquear horario?}]
    Los botones de bloqueo manual son exclusivos de administradores. Para bloquear tu horario, usa \textbf{Ir a Comer} en Asistencia.
\end{description}

% ════════════════════════════════════════════════════════════════
\vspace{2cm}
\begin{center}
  \textcolor{gris}{\rule{0.5\linewidth}{0.4pt}}\\[0.3cm]
  \textcolor{gris}{\small MUYMUY Beauty Studio --- Manual de Empleado v1.1}\\
  \textcolor{primario}{\small Documento de uso interno y confidencial}
\end{center}

\end{document}
